% Work-MAT | 2. rocnik | pripravovana hodina 51
% Návaznost: uživatelské připravy.zip (hodiny 01–20), autorská učebnice na GitHubu.
% Realisticky.cz: didaktická inspirace v posloupnosti a typech úloh; zadání jsou vlastní.
\documentclass[11pt,a4paper]{article}
\input{../../sablony/pripravy/priprava.sty}
\usepackage{tikz}
\nastavPripravu{2. ročník}{Goniometrické funkce}{51}{Tangens a kotangens}
\begin{document}
\blok{Motivace a návaznost}
Navážeme na sinus a kosinus jejich podílem. Ptej se na místa, kde by se dělilo nulou.
\poznamka{Orientační nabídka pro 45--55 minut: krátké zopakování a motivace, nový výklad, tři společné vzorové úlohy, samostatná práce. Vyber podle tempa třídy; zbývající zadání slouží jako rezerva.}
\blok{Výklad}
$\operatorname{tg}\alpha=\frac{\sin\alpha}{\cos\alpha}$ pro $\cos\alpha\ne0$ a $\operatorname{cotg}\alpha=\frac{\cos\alpha}{\sin\alpha}$ pro $\sin\alpha\ne0$.
\blok{Společné příklady}
\begin{prikladReseniBox}[title={Příklad 1}]
Vyřeš, resp. urči: $\operatorname{tg}135^\circ$.
\tcblower
\[\begin{aligned}\frac{\sin135^\circ}{\cos135^\circ}=\frac{\sqrt2/2}{-\sqrt2/2}=-1\end{aligned}\]
\end{prikladReseniBox}
\begin{prikladReseniBox}[title={Příklad 2}]
Vyřeš, resp. urči: $\operatorname{cotg}60^\circ$.
\tcblower
\[\begin{aligned}\frac{\frac{1}{2}}{\frac{\sqrt3}{2}}=\frac{\sqrt3}{3}\end{aligned}\]
\end{prikladReseniBox}
\poznamka{Nauč žáky najít podmínky v každé úloze dříve než dosadí do podílu.}
\blok{Samostatné procvičování}
\begin{zadaniReseniBox}[title={Samostatná práce 1}]
\begin{enumerate}[label=\alph*)]
\item $\cos300^\circ$.
\item $\sin270^\circ$.
\item $\cos180^\circ$.
\item $\operatorname{tg}45^\circ$.
\item $\sin0^\circ$.
\item $\operatorname{cotg}135^\circ$.
\end{enumerate}
\tcblower
\begin{enumerate}[label=\alph*)]
\item $\frac{1}{2}$.
\item $-1$.
\item $-1$.
\item $1$.
\item $0$.
\item $-1$.
\end{enumerate}
\end{zadaniReseniBox}
\blok{Dobrovolné rozšíření}
\begin{zadaniReseniBox}[title={Rozšíření 1}]
\begin{enumerate}[label=\alph*)]
\item Navrhni obdobnou úlohu s jiným úhlem a zkontroluj výsledek graficky..
\end{enumerate}
\tcblower
\begin{enumerate}[label=\alph*)]
\item $Srovnej s jednotkovou kružnicí nebo grafem.$.
\end{enumerate}
\end{zadaniReseniBox}
\blok{Úkol na doma}
Každou hodnotu zdůvodni náčrtkem jednotkové kružnice nebo orientovaným úhlem; u přesných hodnot neuváděj jen desetinné aproximace.
\begin{zadaniReseniBox}[title={Domácí úkol}]
\begin{enumerate}[label=\alph*)]
\item $\sin60^\circ$.
\item $\sin30^\circ$.
\item $\cos60^\circ$.
\item $\sin45^\circ$.
\item $\cos45^\circ$.
\item $\sin150^\circ$.
\item $\cos150^\circ$.
\item $\operatorname{tg}135^\circ$.
\end{enumerate}
\tcblower
\begin{enumerate}[label=\alph*)]
\item $\frac{\sqrt3}{2}$.
\item $\frac{1}{2}$.
\item $\frac{1}{2}$.
\item $\frac{\sqrt2}{2}$.
\item $\frac{\sqrt2}{2}$.
\item $\frac{1}{2}$.
\item $-\frac{\sqrt3}{2}$.
\item $-1$.
\end{enumerate}
\end{zadaniReseniBox}
\poznamka{Domácí úkol obsahuje 8 položek; předpokládaná domácí práce přibližně 15--30 minut. Po větších blocích jsou vědomě promíchané dříve probrané typy.}
\end{document}
